% GENERATED FILE. Edit canonical lessons, then run npm run generate:books.
% canonical-source-hash: fnv1a64:0b717a4f2d96b25b
% canonical-lessons: GU-R18-how-are-you-r4, GU-R18-majaa-r4, GU-R18-no-problem-r4, GU-R18-wellbeing-r4, GU-R18-farewell-r4

\chapter{Fourth Return: Wellbeing and Farewells}
\label{ch:gu-r4-wellbeing-farewells}
\hlchaptermodality{22}

\begin{chapteropening}
\textbf{By the end of this chapter:} \emph{I can retrieve and write the respectful wellbeing exchange and two familiar farewells without a model.}

Five zero-atom lessons close the next exact R4 windows and end with independent farewell production at coherent R3 and R4 intervals.
\end{chapteropening}

\section[Practice]{The respectful wellbeing question returns at R4}

\label{lesson:GU-R18-how-are-you-r4}



\begin{quote}
Ask a new acquaintance \emph{How are you?} without looking.
\end{quote}

\subsection*{Your turn}
{[}YOU SAY: \emph{tũ kem chhe? — tame kem chho?}{]} Choose the respectful question. The pronoun and copula agree as one frame.

\subsection*{Writing — from sound}
Write the three-word question from meaning. Score listening choice and writing separately.

\begin{tcolorbox}[breakable,colback=teal!4,colframe=teal!35!black,title={Before you move on}]
Ask the respectful question once more aloud.
\end{tcolorbox}
\section[Practice]{The wellbeing answer returns at R4}

\label{lesson:GU-R18-majaa-r4}



\begin{quote}
Answer \emph{tame kem chho?} with \emph{I am well} before looking.
\end{quote}

\subsection*{Your turn}
{[}YOU SAY: \emph{tame kem chho? — hũ majāmā chhũ}{]} Choose the answer, then say the question and answer as a pair.

\subsection*{Writing — from sound}
Write the answer from meaning. Listening recognition and written production each have their own score.

\begin{tcolorbox}[breakable,colback=teal!4,colframe=teal!35!black,title={Before you move on}]
What does \emph{majā} literally evoke? (Enjoyment or relish.)
\end{tcolorbox}
\section[Practice]{No problem returns at R4}

\label{lesson:GU-R18-no-problem-r4}



\begin{quote}
Someone says \emph{ābhār}. Give the gracious reply before looking.
\end{quote}

\subsection*{Your turn}
{[}YOU SAY: \emph{ābhār — vāndho nahĩ}{]} Give the reply. Remember that \emph{nahĩ} carries the old negative also heard in \emph{nā}.

\subsection*{Writing — from sound}
Write the reply from meaning. Score the listening choice and the two written words separately.

\begin{tcolorbox}[breakable,colback=teal!4,colframe=teal!35!black,title={Before you move on}]
Say \emph{thanks / no problem} as a two-line exchange.
\end{tcolorbox}
\section[Practice]{The wellbeing exchange returns at R4}

\label{lesson:GU-R18-wellbeing-r4}



\begin{quote}
Picture a question, a positive answer, and a courteous close.
\end{quote}

\subsection*{Your turn}
{[}YOU SAY: \emph{vāndho nahĩ — hũ majāmā chhũ — tame kem chho}{]} Rebuild question, answer, courtesy close.

\subsection*{Writing — from sound}
Write all three lines from meaning. Listening order and written production each need \textbf{3/3}.

\begin{tcolorbox}[breakable,colback=teal!4,colframe=teal!35!black,title={Before you move on}]
Perform the exchange once more aloud.
\end{tcolorbox}
\section[Practice]{We will meet returns at R4}

\label{lesson:GU-R18-farewell-r4}



\begin{quote}
Say \emph{we will meet} before looking.
\end{quote}

\subsection*{Your turn}
{[}YOU SAY: \emph{kāle maḷīshũ — pāchhā maḷīshũ}{]} Give each meaning. \emph{Maḷīshũ} is the shared future verb; \emph{kāle} and \emph{pāchhā} choose the farewell.

\subsection*{Writing — from sound}
Write both farewells from meaning. Score listening meanings and writing separately. This also closes the coherent farewell atoms' R3 windows.

\begin{tcolorbox}[breakable,colback=teal!4,colframe=teal!35!black,title={Before you move on}]
Which verb is shared by both lines? (\textbf{\gu{મળીશું}}, \emph{maḷīshũ}.)
\end{tcolorbox}
